\section{Personaggi del database}
Infine vedremo le principali figure che si interfecciano 
e lavorano col database. \\
Il primo tra tutti è il \textbf{Database Administrator} (\textbf{DBA}),
che definisce lo schema della base di dati, gestisce la sicurezza del database
(quindi decide chi è autorizzato ad accedere a cosa), si occupa della 
manutenzione, delle prestezaioni e dei backup del database e, infine, 
coordina gestisce le esigenze dei vari utenti; per esempio se una azienda A 
vuole accedere agli stessi dati di una azienda B il \textbf{DBA} progetterà
la base di dati in modo tale che entrambe le aziende siano soddisfatte.
Abbiamo inoltre la figura del \textbf{programmatore}, che si occupa
di scrivere applicazioni che interfaccino l'utente col database, in base 
al tipo di utente ci saranno ovviamente applicativi diversi;
per esempio un utente comune non avrà sicuramente accesso a \underline{ogni}
dato del database di una banca. \\
In modo simile gli \textbf{analisti} invece sono utenti ma con conoscenze del
linguaggio SQL, quindi non usano applicativi ma sono loro stessi che "interrogano"
il database con query scritte da loro. \\
Infine gli \textbf{utenti finali} (detti \textbf{naviganti}) sono gli utenti
più comuni, quelli che interagiscono con le applicazioni scritte dai 
\textit{programmatori} già citati prima.